\chapter{Difusión}

El cuidado ambiental es un ámbito de interés público, puesto que los daños que este sufre afectan a toda la población de una comunidad. Así pues, la educación de una población en el cuidado de este es fundamental para que tenga un crecimiento sustentable.

El programa de difusión del club ambiental será una pieza clave para expandir el conocimiento e interés acerca del cuidado ambiental y la importancia que este tiene para la sociedad humana y para la naturaleza.

\subsection{Actividades}

Las actividades que este programa incluirían podrian ser conferencias ,talleres o sesiones de cine. Estás y otras actividades que puedan surgir servirán para dar a conocer el ámbito del cuidado ambiental, acciones que pueden tomar para llevarlo a cabo y la importancia de tal ámbito.

Como se mencionó anteriormente, las actividades de difusión podrían realizarse tanto dentro como fuera de la universidad. En caso de realizarse dentro de esta, el sitio a elegir dependerá del tipo de actividad. Más adelante se especifican los sitios más adecuados por tipo de actividad. En caso de realizarse fuera de la universidad, todas las actividades podrían llevarse a cabo en instituciones educativas (desde primaria hasta universidad) y comunidades civiles (reunir vecinos de una colonia, por ejemplo)

\subsubsection{Conferencias}

Una de las principales actividades de difusión que se podrían llevar a cabo serían conferencias acerca de temas de ambientalismo y tecnologías sustentables. Estas serían organizadas y llevadas a la práctica tanto por integrantes del club ambiental como por alumnos interesados de manera puntual en organizar y/o presentar una conferencia.

Dentro de la universidad, los espacios más adecuados para realizar conferencias podrían ser:

\begin{itemize}
    \item Auditorio de UD-1
    \item Auditorio de UD-2
    \item Sala isóptica
    \item CID
\end{itemize}

\subsubsection{Talleres}

Muchas habilidades y prácticas relacionadas al ambientalismo y la sustentabilidad pueden ser compartidas a las personas a través de talleres. Algunas ideas de talleres podrían ser:

\begin{itemize}
    \item Compostaje en casa y comunitario
    \item Huertos en casa y comunitarios
    \item Preparación casera de alimentos y bebidas
    \item Reparación de bienes (ropa, por ejemplo)
    \item Observación de flora y/o fauna
\end{itemize}

Estos talleres pueden ser llevados a cabo dentro de la universidad en áreas verdes de la universidad y el edificio de tecnologías (“Pecera”).

\subsubsection{Sesiones de cine}

El cine ha sido siempre un medio para comunicar ideas de manera efectiva y entretenida, y existen muchas películas que tratan temas sobre ambientalismo y sustentabilidad las cuales serían adecuadas para compartir con la comunidad estudiantil.

Para llevar a cabo sesiones de cine se podrían reservar los auditorios de la UD-1 y UD-2.

\subsubsection{Actividades recreativas}

Actividades como el ``SpoGomi'' en Japón, han mostrado que se puede combinar el ambientalismo con actividades recreativas y competitivas, por lo que este tipo de actividades también se tendrán en cuenta para su realización.
